\documentclass[11pt,a4paper]{article}
\usepackage[utf8]{inputenc}
\usepackage[T1]{fontenc}
\usepackage[protrusion=true,expansion=false]{microtype}
\usepackage{amsmath}
\usepackage{amssymb}
\usepackage{amsfonts}
\usepackage{amsthm}
\usepackage{graphicx}
\usepackage{cite}
\usepackage{hyperref}
\usepackage{xcolor}
\usepackage{authblk}
\usepackage{algorithm}
\usepackage{algpseudocode}
\usepackage{booktabs}
\usepackage{array}
\usepackage{enumitem}
\usepackage{xurl}

%% Theorem environments
\newtheorem{theorem}{Theorem}[section]
\newtheorem{proposition}[theorem]{Proposition}
\newtheorem{corollary}[theorem]{Corollary}
\newtheorem{definition}[theorem]{Definition}
\newtheorem{remark}[theorem]{Remark}
\newtheorem{conjecture}[theorem]{Conjecture}

\hypersetup{
  colorlinks=true,
  linkcolor=blue!50!black,
  citecolor=blue!50!black,
  urlcolor=blue!50!black,
}

\title{HAML: Hierarchical Attractor Machine Learning\\
with Bidirectional Coupling for Supervised Classification\\[0.5em]
\large A Preliminary Investigation}

\author{Frederick Murat}
\affil{Independent Researcher \\ \texttt{contact via GitHub: ikobootloader}}

\date{May 2026 (v4)}

\begin{document}

\maketitle

\begin{abstract}
We introduce HAML (Hierarchical Attractor Machine Learning), an exploratory supervised
classification framework in which classification emerges from the convergence of a coupled
dynamical system toward stable basins of attraction across multiple hierarchical levels.
Each level operates in a PCA-decomposed subspace of decreasing dimension; levels are coupled
through bottom-up (error-prediction) and top-down (context-prediction) feedback flows
inspired by Predictive Coding.

We present four theoretical contributions:
(1) a dissipative-convergence argument via LaSalle's invariance principle (Proposition~\ref{prop:convergence});
(2) a dimension-independent local convergence-rate result at fixed points under Gaussian
potentials with dimension-scaled kernel initialisation (Proposition~\ref{prop:dimindep});
(3) a hierarchical-decomposition argument giving exponential gains in attractor count
on problems of high topological complexity, with a sub-linear Rademacher bound in depth
(Propositions~\ref{prop:expressivity}--\ref{prop:rademacher});
(4) an exact correspondence with hierarchical Predictive Coding MAP inference under
conditions M1--M3 (Theorem~\ref{thm:map}).

We further introduce a three-phase progressive training strategy to stabilise gradient
learning of coupled dynamics, and a level-aware collapse-recovery mechanism for inter-level
representation collapse.

\textbf{Empirical results (honest summary).}
The empirical study consists of six campaigns; we summarise them with their actual
statistical status rather than headline numbers:
\begin{itemize}
  \item On a MNIST subset (1500 train / 500 test, 5 epochs), gradient training raises test
        accuracy from $45.6\%$ (k-means-only initialisation) to $75.8\%$. On this short
        protocol, bidirectional coupling provides \textbf{no measurable gain} over the
        independent-levels variant ($75.8\%$ vs.\ $75.8\%$), which we discuss explicitly
        rather than hide.
  \item On alternating concentric rings (5 seeds, 42--46), coupled HAML reaches
        $71.20\% \pm 6.43$ vs.\ $67.30\% \pm 4.37$ for the independent variant
        ($\Delta = +3.90$~pts, \emph{not statistically significant} at this sample size).
        The qualitative recovery of non-convex decision boundaries by the coupled model is
        nonetheless visible in the learned geometry.
  \item On noisy moons (single seed), the coupled gain peaks at intermediate noise
        ($+6.27$~pts at noise $0.20$). This is reported as a single-seed observation
        consistent with the theoretical prediction; multi-seed validation is required.
  \item On Fashion-MNIST (784D, 10 classes, 5000 train / 1000 test, single seed,
        10 epochs), test accuracy reaches $80.1\%$. We do \emph{not} claim this is
        competitive with standard baselines (a logistic regression typically reaches
        $\sim$84\% and a small CNN $>90\%$ on this dataset); the figure demonstrates that
        the framework trains stably in 784D rather than that it outperforms standard
        methods.
  \item A controlled \emph{reimplementation} of SA-nODE~\cite{marino2024} (single seed,
        not using the authors' official code) is dominated by HAML across noise levels;
        we emphasise the preliminary nature of this comparison.
\end{itemize}

\textbf{Status of this work.} This is an exploratory preprint from an independent
researcher. The theoretical and architectural contributions are intended as a coherent
proposal; the empirical evidence is preliminary and explicitly under-powered (single
seeds on most real-data benchmarks, no comparison to strong neural baselines, internal
SA-nODE reimplementation). The repository contains the full code and configuration files
required to reproduce every reported number; we welcome external replication.

\textbf{Positioning.} HAML sits at the intersection of four established research
strands --- supervised attractor-based classification, hierarchical Predictive Coding,
hierarchical associative memories, and implicit-depth/equilibrium models. We do not
claim novelty for any individual building block. Our contribution is a specific
instantiation combining (a) supervised classification by basin membership, (b) a
hierarchy of PCA-decomposed subspaces, (c) continuous attractors learned by gradient
descent at every level, and (d) bidirectional PC-style coupling between adjacent levels,
together with theoretical results specialised to this architecture. Detailed
positioning with respect to prior work is provided in Section~\ref{sec:related}.

\textbf{Keywords:} Supervised learning, attractor dynamics, hierarchical classification,
neural ODE, predictive coding, energy-based models, dynamical systems.
\end{abstract}

%% ============================================================
\section{Introduction}
%% ============================================================

Supervised classification traditionally relies on static function approximation through
neural networks or kernel methods. We explore an alternative perspective:
\emph{classification as dynamic convergence}. Rather than computing a forward mapping
$f : \mathcal{X} \to \mathcal{Y}$ in one pass, the model learns to position attractors
(stable equilibrium points) in hierarchical latent subspaces so that unseen samples
converge to basins corresponding to their true class under coupled gradient dynamics.

This approach builds on four established research strands:
\begin{itemize}
  \item \textbf{Dynamical systems theory}: stable manifolds and attractor landscapes provide
        geometrically interpretable decision boundaries.
  \item \textbf{Predictive coding}~\cite{rao1999predictive,friston2005theory}: top-down
        predictions correct bottom-up prediction errors across a generative hierarchy.
  \item \textbf{Energy-based models}: classification emerges from minimisation of a learned
        potential energy.
  \item \textbf{Modern Hopfield networks}~\cite{ramsauer2021hopfield}: continuous attractors
        replace discrete states for greater expressivity.
\end{itemize}

\subsection{Scope and Disclaimer}

We state upfront what this paper does and does not claim, to set expectations correctly:

\textbf{What this paper claims.} (i)~A coherent architectural proposal combining
hierarchical PCA-decomposed state spaces, learned continuous attractors, and bidirectional
Predictive-Coding-style coupling. (ii)~Four theoretical results of the form stated in the
abstract, with their conditions made explicit. (iii)~A practical training recipe
(three-phase progressive coupling, level-aware collapse recovery) that empirically
stabilises gradient-based optimisation of the coupled dynamics. (iv)~Preliminary empirical
evidence that the framework trains stably and produces interpretable attractor geometries.

\textbf{What this paper does \emph{not} claim.} (i)~State-of-the-art results on any
benchmark. (ii)~A definitive empirical advantage of bidirectional coupling: the evidence
is mixed (positive on rings and noisy moons at intermediate noise, neutral on MNIST).
(iii)~Superiority over SA-nODE based on the official code of Marino et al.~\cite{marino2024}
(we use a controlled internal reimplementation). (iv)~Multi-seed certification on real
high-dimensional data: Fashion-MNIST and MNIST results are single-seed for compute
reasons, and we treat them as such.

\subsection{Related Work and Positioning}
\label{sec:related}

HAML sits at the intersection of four research strands. We position our contribution
explicitly with respect to each, listing the closest prior work and the specific
differences with our framework.

\textbf{(i) Attractor-based supervised classification.} The idea of classifying via
convergence to a basin of attraction has a long history. Amit and
Mascaro~\cite{amit2001attractor} described a recurrent perceptron network in which each
class is a pre-learned attractor of the recurrent dynamics, with the visual stimulus
driving the system toward the correct class subset. Dense Associative
Memory~\cite{krotov2016dense} and Modern Hopfield
Networks~\cite{ramsauer2021hopfield} re-energised this line with continuous attractors and
exponential storage capacity. Most directly,
SA-nODE~\cite{marino2024} introduces a continuous-time supervised classifier with
pre-planted stable attractors in a single ambient space, trained end-to-end as a Neural
ODE~\cite{chen2018neural}. HAML differs from these by (a) operating in a
\emph{hierarchy} of subspaces rather than a single ambient space, (b) learning
\emph{continuous} attractor positions, widths, and repulsion radii by gradient descent
rather than planting them analytically, and (c) introducing explicit inter-level
coupling.

\textbf{(ii) Hierarchical Predictive Coding.} Predictive Coding (PC) posits a
hierarchical generative model with top-down predictions and bottom-up prediction
errors~\cite{rao1999predictive,friston2005theory}.
Spratling~\cite{spratling2017hierarchical} demonstrated that hierarchical PC can perform
visual object recognition; Han et al.~\cite{han2018deep} introduced a deep bidirectional
PC network with local recurrent processing whose internal representations \emph{converge
over time} during classification on SVHN, CIFAR, and ImageNet --- structurally the
closest prior work to HAML's bidirectional coupling. The mathematical relationship
between PC and backpropagation has been studied in
depth~\cite{whittington2017approximation,song2020can,salvatori2021exactbp}, with
results showing that PC networks can approximate or exactly replicate backpropagation
under specific conditions. HAML differs from this line by (a) framing the dynamics as a
continuous-time dissipative system with explicit attractor structure rather than as a
discrete recurrent network, (b) using PCA-based subspace decomposition between levels,
and (c) reading out classification directly from basin membership rather than from
top-layer activity.

\textbf{(iii) Hierarchical associative memories.} Krotov~\cite{krotov2021hierarchical}
introduced a fully recurrent associative memory model with an arbitrary number of layers
and a Lyapunov energy function that decreases along the dynamics. Salvatori et
al.~\cite{salvatori2021associative} realised associative memories via a hierarchical
generative network trained by predictive coding, with retrieval framed as inference.
Both are very close to HAML's hierarchical energy/attractor framing but target
auto-associative memory and retrieval rather than supervised classification with
class-specific basins.

\textbf{(iv) Implicit-depth and equilibrium models.} Deep Equilibrium Models
(DEQ)~\cite{bai2019deep} compute outputs as fixed points of an implicit
transformation, with $O(1)$ memory via implicit differentiation; Multiscale Deep
Equilibrium Models (MDEQ)~\cite{bai2020multiscale} extend this to hierarchical
multi-resolution feature maps and achieve competitive results on ImageNet and
Cityscapes. Monotone Operator Equilibrium Networks
(monDEQ)~\cite{winston2020monotone} provide existence and uniqueness guarantees under
monotonicity constraints. Generalisation bounds for neural ODEs have been derived via
the Chen--Fliess series~\cite{hanson2024rademacher} and via path-norm
arguments~\cite{marion2023generalization}. These works share with HAML the use of an
implicit/continuous dynamics whose endpoint defines the output, but do not use a
hierarchical attractor structure for class decision or a Predictive-Coding-style
bidirectional coupling.

\textbf{Where HAML sits.} We do not claim that any individual building block
(attractors, hierarchical PC, energy minimisation, fixed-point computation, dimension-
scaled kernels) is new. To our knowledge, we did not find a published model that
combines simultaneously: (a) a supervised-classification objective with read-out by
basin membership, (b) a hierarchy of subspaces obtained by PCA decomposition,
(c) continuous attractors learned by gradient descent at every level, and
(d) a bidirectional PC-style coupling between adjacent levels. HAML's contribution is
this specific instantiation, together with the theoretical results that specialise to
it. We treat the closeness of each building block to prior work explicitly when
discussing limitations (Section~\ref{sec:discussion:limitations}).

Table~\ref{tab:comparison_sanode} summarises the architectural differences with
SA-nODE, the closest prior work in the
\emph{supervised-attractor-classification-via-ODE} family.

\begin{table}[h]
\centering
\footnotesize
\setlength{\tabcolsep}{5pt}
\begin{tabular}{lll}
\toprule
\textbf{Dimension} & \textbf{SA-nODE} & \textbf{HAML} \\
\midrule
State space             & Full input dim.\               & Hierarchy of decreasing subspaces \\
Attractors              & Binary eigenvectors $\pm a$    & Continuous, gradient-optimised \\
Potential               & Fixed double-well              & Parametric Gaussian mixture \\
Structure               & Single dynamic level           & $L$ bidirectionally coupled levels \\
Coupling                & None                           & Bottom-up + top-down \\
Local conv.\ rate at FP & Dimension-dependent            & Dim.-independent (Prop.~\ref{prop:dimindep}) \\
Theoretical anchor      & Statistical mechanics          & Dyn.\ systems + Predictive Coding \\
PC/MAP equivalence      & Not established                & Exact under M1--M3 (Thm.~\ref{thm:map}) \\
\bottomrule
\end{tabular}
\caption{Architectural comparison with SA-nODE~\cite{marino2024}.}
\label{tab:comparison_sanode}
\end{table}

\subsection{Contributions}

\begin{enumerate}
  \item \textbf{Architecture}: a hierarchical bidirectional attractor system with
        gradient-learned continuous attractors for supervised classification.
  \item \textbf{Four theoretical results}: dissipative convergence (Proposition~\ref{prop:convergence}),
        dimension-independent local rate at fixed points (Proposition~\ref{prop:dimindep}),
        hierarchical expressivity gain (Propositions~\ref{prop:expressivity}--\ref{prop:rademacher}),
        and exact MAP/Predictive Coding correspondence under M1--M3 (Theorem~\ref{thm:map}).
  \item \textbf{Preliminary empirical study}: 6 campaigns across synthetic and real
        datasets, with multi-seed runs where compute allowed.
  \item \textbf{Practical mechanisms}: three-phase progressive training and level-aware
        collapse recovery.
  \item \textbf{Open reproducibility}: full source code at
        \url{https://github.com/ikobootloader/HIERARCHICAL-ATTRACTOR-MACHINE-LEARNING}.
\end{enumerate}

%% ============================================================
\section{Model Architecture and Dynamics}
%% ============================================================

\subsection{Hierarchical State Space}

Let $\mathcal{X} \subset \mathbb{R}^d$ be the input space. We define a tower of linear
projections:
\begin{equation}
  \mathcal{X} = \mathcal{H}_0
  \xrightarrow{\pi_1} \mathcal{H}_1
  \xrightarrow{\pi_2} \cdots
  \xrightarrow{\pi_L} \mathcal{H}_L
\end{equation}
where $\mathcal{H}_l \subset \mathbb{R}^{d_l}$ with $d_0 > d_1 > \cdots > d_L$.

Projections $\pi_l : \mathcal{H}_{l-1} \to \mathcal{H}_l$ are initialised by PCA (capturing
maximum variance) and are optionally fine-tuned during training. Their Moore--Penrose
pseudo-inverses $\pi_l^\dagger : \mathcal{H}_l \to \mathcal{H}_{l-1}$ serve as upward
decoders.

\begin{remark}[Contraction property]
PCA projections normalised to unit spectral norm satisfy
$\|\Pi_l\|_2 \le \sqrt{d_{l+1}/d_l} < 1$.
This contraction is used in the convergence analysis (Section~\ref{sec:theory:convergence}).
\end{remark}

Initial conditions: $\mathbf{x}^{(l)}(0) = \pi_l \circ \cdots \circ \pi_1(\mathbf{x}_{\mathrm{input}})$.

\subsection{Attractors and Energy Potential}

At level $l$, we maintain $K \times M$ learnable attractors
$\{\boldsymbol{\mu}^{(l)}_{c,m}\}_{c \in [K], m \in [M]}$
with associated parameters $(\sigma^{(l)}_{c,m}, \rho^{(l)}_{c,m}, w^{(l)}_{c,m})$.
The energy potential is a parametric attractor--repulsor mixture:

\begin{equation}
  E^{(l)}(\mathbf{x}) = -\sum_{c=1}^{K} \sum_{m=1}^{M} w_{c,m}^{(l)}
    \exp\!\left(-\frac{\|\mathbf{x} - \boldsymbol{\mu}_{c,m}^{(l)}\|^2}
                      {2(\sigma_{c,m}^{(l)})^2}\right)
  + \lambda \sum_{c \ne c'} \sum_{m,m'} w_{c',m'}^{(l)}
    \exp\!\left(-\frac{\|\mathbf{x} - \boldsymbol{\mu}_{c',m'}^{(l)}\|^2}
                      {2(\rho_{c',m'}^{(l)})^2}\right)
\label{eq:energy}
\end{equation}

\textbf{Bayesian interpretation} (anticipating Section~\ref{sec:theory:bayes}):
without the repulsive terms ($\lambda=0$), $-E^{(l)}$ is proportional to the log of a
Gaussian mixture, i.e.\ the log-prior $\log p(\mathbf{x}^{(l)})$ in the implicit generative
model. Repulsive terms correspond to a discriminative contrastive
prior~\cite{larochelle2008}.

\subsection{Coupled Dynamics}

At each level $l \in \{0, \ldots, L\}$, the state evolves under the coupled ODE:
\begin{equation}
  \dot{\mathbf{x}}^{(l)}(t) =
  \mathbf{F}^{(l)}_{\mathrm{intra}}(\mathbf{x}^{(l)})
  + \mathbf{F}^{(l)}_{\mathrm{bu}}(\mathbf{x}^{(l)}, \mathbf{x}^{(l-1)})
  + \mathbf{F}^{(l)}_{\mathrm{td}}(\mathbf{x}^{(l)}, \mathbf{x}^{(l+1)})
  - \gamma \mathbf{x}^{(l)}
\label{eq:dynamics}
\end{equation}
where $\gamma > 0$ is a viscous damping coefficient ensuring energy dissipation.

\medskip
\textbf{Intra-level force (local attractors and repulsors):}
\begin{equation}
  \mathbf{F}^{(l)}_{\mathrm{intra}} = -\nabla_{\mathbf{x}} E^{(l)}(\mathbf{x}^{(l)})
\end{equation}

\textbf{Bottom-up error-correction force:}
\begin{equation}
  \mathbf{F}^{(l)}_{\mathrm{bu}} =
  \alpha_{\mathrm{bu}} \cdot \pi_l\!\left(
    \mathbf{x}^{(l-1)} - \pi_l^\dagger(\mathbf{x}^{(l)})
  \right)
\end{equation}
This is the \emph{prediction error} signal of Predictive Coding~\cite{rao1999predictive}:
the difference between the actual lower-level state and the reconstruction predicted by
the current level.

\textbf{Top-down prediction force:}
\begin{equation}
  \mathbf{F}^{(l)}_{\mathrm{td}} =
  \alpha_{\mathrm{td}} \cdot \left(
    \pi_l^\dagger\!\left(\mathbf{x}^{(l+1)}\right) - \mathbf{x}^{(l)}
  \right)
\end{equation}
This is the recall force toward the representation predicted by the higher level --- the
\emph{top-down generation} in the Predictive Coding sense.

\textbf{Boundary conditions}: $\mathbf{F}^{(0)}_{\mathrm{bu}} = 0$,
$\mathbf{F}^{(L)}_{\mathrm{td}} = 0$.

\subsection{Prediction}

After integration until steady state (criterion: $\|\dot{\mathbf{x}}^{(l)}\| < \epsilon\;\forall l$):
\begin{equation}
  \hat{c} = \operatorname*{arg\,max}_c \sum_{l=0}^{L} \beta_l \sum_{m=1}^{M}
    \exp\!\left(-\frac{\|\mathbf{x}^{(l)}(T) - \boldsymbol{\mu}^{(l)}_{c,m}\|^2}
                      {2(\sigma^{(l)}_{c,m})^2}\right)
\end{equation}
where $\beta_l$ are per-level weights (default: $\beta_l = 2^l$, learned or fixed).

\subsection{Critical Design: Dimension-Dependent Kernel Initialisation (Condition I1)}

\textbf{Problem}: in high dimension $d$, Euclidean distances concentrate around
$\mathbb{E}[\|\mathbf{x}-\mathbf{y}\|] \approx \sqrt{d}\cdot\sigma_{\mathrm{data}}$.
With a fixed $\sigma_{\mathrm{kernel}} \sim 1.0$, Gaussian kernels saturate toward~0 for
$d \gtrsim 20$, forcing $\nabla E \to 0$ and halting learning.

\textbf{Solution (I1)}: initialise with dimension scaling:
\begin{equation}
  \sigma^{(l)}_{\mathrm{init}} = \sqrt{d_l} \cdot \sigma^{(l)}_{\mathrm{data}}
\end{equation}

\textbf{Empirical impact}:

\begin{table}[h]
\centering
\begin{tabular}{rrrr}
\toprule
\textbf{Dimension} & \textbf{Pre-fix kernel} & \textbf{Post-fix kernel} & \textbf{Gradient magnitude} \\
\midrule
10  & 0.95     & 0.88 & 0.24 \\
100 & 0.03     & 0.67 & 0.18 \\
784 & $<0.01$  & 0.52 & 0.15 \\
\bottomrule
\end{tabular}
\caption{Kernel values and gradient magnitude by dimension (typical single run).
Without~I1, the framework fails to train beyond $\sim$20D in our experiments.}
\label{tab:kernelfix}
\end{table}

%% ============================================================
\section{Learning via Gradient Descent}
%% ============================================================

\subsection{Three-Phase Progressive Training}

To stabilise learning of coupled dynamics, we use a three-phase schedule:

\textbf{Phase~1 (epochs $1$--$E_1$): Independent learning.}
Set $\alpha_{\mathrm{bu}} = \alpha_{\mathrm{td}} = 0$. Each level self-organises
independently. Ensures conditions C1 and C3 (see Section~\ref{sec:theory:convergence}) are
empirically satisfied before coupling is introduced.

\textbf{Phase~2 (epochs $E_1+1$--$E_1+E_2$): Progressive coupling.}
\begin{equation}
  \alpha_{\mathrm{bu}}(t) = \alpha_{\mathrm{bu}}^{\max} \cdot \frac{t - E_1}{E_2},
  \quad
  \alpha_{\mathrm{td}}(t) = \alpha_{\mathrm{td}}^{\max} \cdot \frac{t - E_1}{E_2}
\end{equation}

\textbf{Phase~3 (epochs $E_1+E_2+1$--$E_{\mathrm{total}}$): Full coupling.}
Fix $\alpha_{\mathrm{bu}} = \alpha_{\mathrm{bu}}^{\max}$,
$\alpha_{\mathrm{td}} = \alpha_{\mathrm{td}}^{\max}$.

\subsection{Composite Loss Function}

\begin{equation}
  \mathcal{L} = \mathcal{L}_{\mathrm{CE}} + \mu_1 \mathcal{L}_{\mathrm{sep}} + \mu_2 \mathcal{L}_{\mathrm{dyn}}
\label{eq:loss}
\end{equation}

\textbf{Cross-entropy term} (multi-level softmax):
\begin{equation}
  \mathcal{L}_{\mathrm{CE}} = -\frac{1}{N} \sum_{i=1}^{N} \log P(y_i \mid \{\mathbf{x}_i^{(l)}(T)\}_l)
\end{equation}

\textbf{Basin separation term}:
\begin{equation}
  \mathcal{L}_{\mathrm{sep}} = \sum_{l} \sum_{c \ne c'} \sum_{m,m'}
    \exp\!\left(-\frac{\|\boldsymbol{\mu}^{(l)}_{c,m} - \boldsymbol{\mu}^{(l)}_{c',m'}\|^2}
                       {2\sigma^2_{\min}}\right)
\end{equation}

\textbf{Dynamical regularisation}:
\begin{equation}
  \mathcal{L}_{\mathrm{dyn}} = \frac{1}{T}\int_0^T \sum_l \|\dot{\mathbf{x}}^{(l)}(t)\|^2\, \mathrm{d}t
\end{equation}

Gradients are computed via RK4 unrolling (for short $T$) or the adjoint
method~\cite{chen2018neural} via \texttt{torchdiffeq} (for long $T$, memory-efficient).

Attractor positions $\boldsymbol{\mu}^{(l)}_{c,m}$ are initialised by per-class $k$-means
on the projected training data, providing an initial configuration that empirically
satisfies the local-dominance assumption (C1) on simple problems.

\subsection{Level-Aware Collapse Recovery}
\label{sec:p6}

\begin{proposition}[Level-Aware Recovery, informal mechanism]
\label{prop:p6}
Define a level $l^*$ as ``collapsed'' when its single-level accuracy
$|\mathrm{acc}_{l^*} - 1/K| \le \varepsilon$ for at least \texttt{patience}
consecutive epochs during Phase~3, while global train accuracy is moderate
(anti-false-positive filter). When detected, the following is triggered:
\begin{enumerate}
  \item Re-anchor attractors of level $l^*$ via supervised prototype initialisation in
        its latent space.
  \item Temporarily attenuate top-down forces with cooldown ($\alpha_{\mathrm{td}}$ scaled
        by $\tau<1$ for \texttt{td\_cooldown\_epochs}).
\end{enumerate}
\end{proposition}

This is a heuristic mechanism, not a theorem. Empirically on the concentric-rings
benchmark (Campaign~E8, 5 seeds 42--46), it provides a mean accuracy gain of $+2.2$~pts
and a worst-seed gain of $+4$~pts. We treat this as engineering scaffolding, not as a
contribution at the same epistemic level as the four theoretical propositions.

%% ============================================================
\section{Theoretical Analysis}
\label{sec:theory}
%% ============================================================

\subsection{Convergence Guarantees}
\label{sec:theory:convergence}

In the overdamped regime (large $\gamma$, negligible inertia), the system reduces to a pure
gradient flow:
\begin{equation}
  \dot{\mathbf{X}} = -\frac{1}{\gamma}\nabla_{\mathbf{X}} \mathcal{E}(\mathbf{X})
\end{equation}
where $\mathbf{X} = (\mathbf{x}^{(0)}, \ldots, \mathbf{x}^{(L)})$ and the total energy is:
\begin{equation}
  \mathcal{E}(\mathbf{X}) = \sum_l E^{(l)}(\mathbf{x}^{(l)})
  + \frac{\alpha_{\mathrm{bu}} + \alpha_{\mathrm{td}}}{2}
    \sum_l \|\mathbf{x}^{(l)} - \pi_l^\dagger(\mathbf{x}^{(l+1)})\|^2
\label{eq:lyapunov}
\end{equation}

\begin{proposition}[Dissipative Convergence]
\label{prop:convergence}
$\mathcal{E}$ is bounded below by $-\sum_l \sum_{c,m} w^{(l)}_{c,m}$. Along any trajectory,
$\dot{\mathcal{E}} = -\frac{1}{\gamma}\|\nabla \mathcal{E}\|^2 \le 0$.
By LaSalle's invariance principle, every bounded trajectory converges to the largest
invariant subset of $\{\mathbf{X} : \nabla \mathcal{E}(\mathbf{X}) = 0\}$. The system can
neither diverge nor oscillate indefinitely.
\end{proposition}

\begin{remark}[On the strength of Proposition~\ref{prop:convergence}]
Proposition~\ref{prop:convergence} establishes convergence to \emph{some} critical point
of~$\mathcal{E}$. It does not by itself guarantee convergence to a \emph{minimum},
let alone to the \emph{correct-class} minimum. The following result formalises sufficient
conditions for the latter and is, we acknowledge, largely tautological in flavour: it
states what every gradient method needs.
\end{remark}

\begin{proposition}[Sufficient conditions for correct-basin convergence]
\label{prop:correct_basin}
Assume:
\begin{description}[leftmargin=2.5em]
  \item[C1 -- Local dominance]: at each level, the true-class attractor is the nearest
        attractor to the trajectory's neighbourhood in Euclidean distance.
  \item[C2 -- Sufficient coupling]:
        $\alpha_{\mathrm{bu}} + \alpha_{\mathrm{td}} > \alpha_{\min}$,
        a problem-dependent threshold determined by spectral analysis of the linearised
        system around fixed points.
  \item[C3 -- Basin separation]: $\mathcal{L}_{\mathrm{sep}}$ enforces a positive margin
        between attractors of different classes.
\end{description}
Then any trajectory whose initial condition lies in the joint true-class basin converges
to the true-class fixed point at each level.
\end{proposition}

\begin{remark}[Honest framing]
Proposition~\ref{prop:correct_basin} is essentially: \emph{if} the initial condition lies
in the correct basin (C1) \emph{and} the basins are separated (C3) \emph{and} the coupling
is well-tuned (C2), \emph{then} gradient flow converges in that basin. We state it
explicitly because (a) C1--C3 are precisely the conditions we monitor and enforce via
$k$-means initialisation, $\mathcal{L}_{\mathrm{sep}}$, and the three-phase schedule, and
(b) the architectural value of HAML lies in making C1 hold at \emph{higher} levels even
when it fails at level~0, via top-down guidance (Proposition~\ref{prop:guidance}). The
result is not deep in itself, but it makes the role of each design choice explicit.
\end{remark}

\begin{proposition}[Hierarchical Guidance]
\label{prop:guidance}
Under coupling, the linearised dynamics around a fixed point have a block-tridiagonal
Hessian (Eq.~\ref{eq:hessian}) whose minimum eigenvalue is bounded below by
$\min_l \lambda_{\min}(\mathbf{H}^{(l)}) + (\alpha_{\mathrm{bu}}+\alpha_{\mathrm{td}}) - 2\alpha_{\mathrm{td}}\max_l\|\Pi_l\|_2$.
For PCA projections with $\|\Pi_l\|_2 \le 1$ and moderate coupling, this lower bound is
strictly positive, ensuring exponential local convergence at every level once the
trajectory enters the basin of the higher level.
\end{proposition}

The Hessian of $\mathcal{E}$ at a fixed point has block-tridiagonal structure:
\begin{equation}
  \mathbf{H} = \begin{pmatrix}
    \mathbf{H}^{(0)} + \mathbf{C}^{(0)} & -\alpha_{\mathrm{td}}\Pi_1^T & 0 & \cdots \\
    -\alpha_{\mathrm{td}}\Pi_1 & \mathbf{H}^{(1)} + \mathbf{C}^{(1)} & -\alpha_{\mathrm{td}}\Pi_2^T & \cdots \\
    0 & -\alpha_{\mathrm{td}}\Pi_2 & \ddots & \vdots \\
    \vdots & \vdots & \vdots & \ddots
  \end{pmatrix}
\label{eq:hessian}
\end{equation}
with $\mathbf{C}^{(l)} = (\alpha_{\mathrm{bu}} + \alpha_{\mathrm{td}})\mathbf{I}_{d_l}$.
By Weyl's interlacing inequality~\cite{weyl1912}, the coupling matrices $\mathbf{C}^{(l)}$
add a positive shift to the eigenvalues of the block-diagonal part; the off-diagonal
$-\alpha_{\mathrm{td}}\Pi_l$ blocks act as a perturbation bounded by
$\alpha_{\mathrm{td}}\|\Pi_l\|_2$. In the moderate-coupling regime
$\alpha_{\mathrm{td}} \|\Pi_l\|_2 < \alpha_{\mathrm{bu}}+\alpha_{\mathrm{td}}$, the
coupling strictly accelerates local convergence.

\subsection{Dimension-Independent Local Convergence Rate}
\label{sec:theory:dimindep}

\begin{proposition}[Dimension-Independent Local Rate]
\label{prop:dimindep}
At a fixed point $\mathbf{x}^{(l)*} = \boldsymbol{\mu}^{(l)}_{c,m}$, the intra-level
Hessian block of $E^{(l)}$ in Eq.~\ref{eq:energy} reduces to:
\begin{equation}
  \mathbf{H}^{(l)}\big|_{\boldsymbol{\mu}} = \kappa^{(l)} \cdot \mathbf{I}_{d_l}
\end{equation}
where
\[
\kappa^{(l)} = \frac{w^{(l)}_{c,m}}{(\sigma^{(l)}_{c,m})^2}
  - \lambda \sum_{c' \ne c, m'} \frac{w^{(l)}_{c',m'}}{(\rho^{(l)}_{c',m'})^2}
    \exp\!\left(-\frac{\|\boldsymbol{\mu}^{(l)}_{c,m} - \boldsymbol{\mu}^{(l)}_{c',m'}\|^2}{2(\rho^{(l)}_{c',m'})^2}\right)
\]
is a \textbf{scalar independent of~$d_l$}. The local linear convergence time
$T \sim \gamma/\kappa^{(l)}$ does not grow with the level dimension.
\end{proposition}

\begin{remark}[Scope of Proposition~\ref{prop:dimindep}]
This is a statement about the \emph{local} linear rate at a fixed point. It does not
address (a) the time to \emph{enter} the basin, which depends on the energy landscape
geometry away from the fixed point, nor (b) the global convergence time, which can still
be dimension-sensitive. The result establishes that, \emph{conditional on having reached
the linear regime near an attractor}, the residual relaxation speed is set by a scalar
that does not scale with~$d_l$. Combined with condition~I1
($\sigma_{\mathrm{init}} \propto \sqrt{d_l}$), this explains the empirical observation
that the framework continues to train at 784D where naive constant-$\sigma$ kernels fail
(Table~\ref{tab:kernelfix}).
\end{remark}

\begin{corollary}
\label{cor:depth}
Adding hierarchical levels does not, in itself, slow down local convergence at fixed
points. By the eigenvalue interlacing analysis of Eq.~\ref{eq:hessian},
$\lambda_{\min}(\mathbf{H})$ remains bounded below by a quantity independent of~$L$,
provided $\alpha_{\mathrm{td}}\max_l\|\Pi_l\|_2$ does not grow with~$L$.
\end{corollary}

\subsection{Expressivity of the Hierarchy}
\label{sec:theory:expressivity}

\begin{definition}[Topological complexity]
The topological complexity $\kappa(\mathcal{P})$ of a binary classification problem
$\mathcal{P}$ is the minimum number of connected components of the boundary of a Bayes-optimal
decision region.
\end{definition}

\begin{proposition}[Hierarchical Expressivity Gain]
\label{prop:expressivity}
A flat ($L=1$) attractor model with one Gaussian attractor per simply-connected basin
requires at least $\kappa(\mathcal{P})$ attractors to represent a problem of topological
complexity~$\kappa(\mathcal{P})$. A hierarchical model with $L$ levels and $M$ attractors
per level can represent up to $O(M^L)$ distinct joint configurations $(c,m_0,\ldots,m_L)$
with only $O(ML)$ parameters per class. On problems whose Bayes-optimal decision boundary
factorises hierarchically (i.e.\ where a global discrete choice at level~$L$ selects a
local geometry at level~0), the parameter count gap between flat and hierarchical is
\emph{exponential in~$L$}.
\end{proposition}

\begin{remark}[Caveats]
Proposition~\ref{prop:expressivity} is an \emph{existence}-style statement: it bounds
the worst-case advantage of the hierarchy on a specific structural family of problems
(hierarchically factorisable boundaries). It does not say that hierarchical models always
beat flat ones --- on problems without such factorisation, the gain may be zero. The
concentric-rings problem (Section~\ref{sec:emp:rings}) is a constructive example where the
factorisation holds: radial structure factorises out at level~1, angular refinement
remains at level~0.
\end{remark}

\textbf{Analytical example --- alternating concentric rings}: classes defined by
$r \in [0,1] \cup [2,3]$ (class~1) and $r \in [1,2] \cup [3,4]$ (class~2). With $L=1$
in the ambient 2D space, one needs $\Omega(1/\sigma)$ attractors per ring to cover its
angular extent. With $L=2$, level~1 (radial projection, $d_1=1$) resolves the global
annular structure with 2 attractors per class; the top-down signal then constrains
level~0 to refine locally.

\begin{conjecture}[Necessity of Bidirectional Coupling]
\label{prop:bidi_necessary}
We conjecture that there exist classification problems representable by a hierarchical
bidirectional system with $O(ML)$ attractors that cannot be represented by $L$ independent
attractor levels with the same total attractor count. We do not provide a proof here. The
empirical signature consistent with this conjecture is the qualitative recovery of
non-convex decision boundaries by the coupled model on the concentric-rings benchmark
(Section~\ref{sec:emp:rings}); but mean accuracy alone does not separate the two models
statistically on our current sample size, so this remains conjectural.
\end{conjecture}

\begin{proposition}[Rademacher Complexity Upper Bound]
\label{prop:rademacher}
For the coupled $L$-level classifier defined by Eq.~\ref{eq:dynamics} with $K$ classes,
$M$ attractors per level, and $n$ training samples, an upper bound on the empirical
Rademacher complexity of the induced function class $\mathcal{F}_L$ takes the form:
\begin{equation}
  \mathcal{R}_n(\mathcal{F}_L) \le c \sqrt{\frac{KML}{n}}
    \cdot \prod_{l=1}^{L}(1 + \alpha_l\|\Pi_l\|_F)
\end{equation}
for a universal constant $c$ depending on the Lipschitz constants of the integrated flow.
The dependence on depth is $O(\sqrt{L})$ inside the square root and at most polynomial via
the product term (which is bounded since $\|\Pi_l\|_F \le \sqrt{d_l}$ and PCA spectral
norm is bounded).
\end{proposition}

\begin{remark}[Relation to prior generalisation results for implicit/continuous models]
The existence of Rademacher complexity bounds for neural ODEs and implicit-depth networks
is not new: Hanson and Raginsky~\cite{hanson2024rademacher} derive bounds via the
Chen--Fliess series for neural ODEs, and Marion~\cite{marion2023generalization} provides
generalisation bounds for neural ODEs and deep residual networks under various path-norm
assumptions. The specific structure we exploit here is the
\emph{hierarchical-product form} of the bound: the $O(\sqrt{ML})$ scaling in the number
of attractors and depth is a direct consequence of the union-bound argument over levels,
and the product $\prod_l (1 + \alpha_l\|\Pi_l\|_F)$ reflects the Lipschitz contraction
afforded by PCA projections (cf.\ Remark on contraction in Section~\ref{sec:theory:dimindep}).
We do \emph{not} claim that this bound is tighter than existing bounds applied to the
same architecture; we claim that it is the natural specialisation of the general theory to
HAML's specific structure. Closing the gap with the tightest known bounds for related
implicit models is open work.
\end{remark}

\subsection{Predictive Coding Correspondence and Bayesian Grounding}
\label{sec:theory:bayes}

\subsubsection{The Predictive Coding Generative Model}

Predictive Coding~\cite{rao1999predictive,friston2005theory} posits a hierarchical
generative model:
\begin{equation}
  p(\mathbf{x}^{(0:L)}) = p(\mathbf{x}^{(L)}) \prod_{l=0}^{L-1}
    p(\mathbf{x}^{(l)} \mid \mathbf{x}^{(l+1)})
\end{equation}
In the point-mean-field approximation, the variational free energy reduces to:
\begin{equation}
  \mathcal{F}\big|_{\mathrm{pt}} =
    \sum_l \frac{1}{2}\|\boldsymbol{\varepsilon}^{(l)}\|^2_{\Sigma^{(l)}}
    - \log p(\mathbf{x}^{(L)})
\end{equation}
where $\boldsymbol{\varepsilon}^{(l)} = \mathbf{x}^{(l)} - g^{(l+1)}(\mathbf{x}^{(l+1)})$
are prediction errors. The inference dynamics are:
\begin{equation}
  \dot{\mathbf{x}}^{(l)} = -\frac{\partial \mathcal{F}}{\partial \mathbf{x}^{(l)}}
    = -\boldsymbol{\varepsilon}^{(l)}
    + \left(\frac{\partial g^{(l+1)}}{\partial \mathbf{x}^{(l)}}\right)^T
      \boldsymbol{\varepsilon}^{(l+1)}
\end{equation}

This is the structural form of HAML's bidirectional coupling.

\subsubsection{Term-by-Term Identification}

With linear generators $g^{(l)}(\mathbf{x}^{(l)}) = \Pi_l^\dagger \mathbf{x}^{(l)}$ and
isotropic covariances $\Sigma^{(l)} = \frac{1}{\alpha_{\mathrm{bu}}+\alpha_{\mathrm{td}}}\mathbf{I}$:
\begin{equation}
  \mathcal{F}\big|_{\mathrm{linear}} =
    \frac{\alpha_{\mathrm{bu}}+\alpha_{\mathrm{td}}}{2}
    \sum_l \|\mathbf{x}^{(l)} - \Pi_l^\dagger \mathbf{x}^{(l+1)}\|^2
    - \log p(\mathbf{x}^{(L)})
\end{equation}
The coupling terms are identical to those in $\mathcal{E}$ (Eq.~\ref{eq:lyapunov}).
Identifying $E^{(l)}$ with $-\log p(\mathbf{x}^{(l)})$:
\begin{equation}
  -E^{(l)}(\mathbf{x})\big|_{\lambda=0} =
    \log \sum_{c,m} w^{(l)}_{c,m}\,
    \mathcal{N}(\mathbf{x} \mid \boldsymbol{\mu}^{(l)}_{c,m}, (\sigma^{(l)}_{c,m})^2\mathbf{I})
    + \text{const}
\end{equation}

\begin{theorem}[MAP Equivalence under M1--M3]
\label{thm:map}
Under conditions:
\begin{description}[leftmargin=2.5em]
  \item[M1 -- Linear generators]: $g^{(l)}(\mathbf{x}^{(l)}) = \Pi_l^\dagger \mathbf{x}^{(l)}$
        (satisfied by construction in our architecture).
  \item[M2 -- Isotropic covariances]: $\Sigma^{(l)} = \frac{1}{\alpha_{\mathrm{bu}}+\alpha_{\mathrm{td}}}\mathbf{I}$
        (satisfied when $\alpha_{\mathrm{bu}}$ and $\alpha_{\mathrm{td}}$ are level-independent;
        relaxable to per-level scalings).
  \item[M3 -- No repulsion]: $\lambda = 0$ in Eq.~\ref{eq:energy}.
\end{description}
With these three conditions, $\mathcal{E} \equiv \mathcal{F}|_{\mathrm{pt}}$ as functionals
of $\mathbf{X}$ (up to an additive constant), and HAML's gradient flow exactly performs
\emph{MAP inference} in a hierarchical generative model with Gaussian mixture prior at
each level and Gaussian conditional likelihoods.

With $\lambda > 0$, $\mathcal{E} = \mathcal{F}_{\mathrm{PC}} + \lambda\,\mathcal{F}_{\mathrm{discrim}}$,
where $\mathcal{F}_{\mathrm{discrim}}$ corresponds to a contrastive discriminative
prior~\cite{larochelle2008}. This hybrid is a discriminative augmentation of standard PC
inference.
\end{theorem}

\begin{remark}[Relation to prior PC--inference and PC--backprop correspondences]
\label{rem:map_priorwork}
The correspondence between Predictive Coding and variational inference in hierarchical
Gaussian generative models is well established: it traces back to Rao and
Ballard~\cite{rao1999predictive} and was formalised in Friston's free-energy
framework~\cite{friston2005theory}. The relationship between PC and backpropagation
has been the subject of substantial recent work: Whittington and
Bogacz~\cite{whittington2017approximation} showed that PC networks approximate the
error backpropagation algorithm under specific conditions; Song et
al.~\cite{song2020can} established exact equivalence to BP on multilayer perceptrons
under stronger assumptions; Salvatori et al.~\cite{salvatori2021exactbp} extended this
to arbitrary computation graphs.

Theorem~\ref{thm:map} does \emph{not} claim novelty for the PC--MAP correspondence in
general. The result is a \emph{specialisation} of these established correspondences to
HAML's specific generative model: Gaussian-mixture priors at each level (rather than
Gaussian priors), PCA-projection generators (rather than learned MLPs), and isotropic
inverse covariances tied to the coupling coefficients. We make conditions M1--M3
explicit precisely to highlight that the correspondence is conditional on these
specific structural choices. The genuinely novel element is the integration of this
correspondence with a supervised attractor read-out and a hierarchical Gaussian-mixture
prior, not the abstract fact that gradient descent on $\mathcal{E}$ equals
free-energy minimisation.
\end{remark}

The implicit generative model is:
\begin{align}
  p(\mathbf{x}^{(l)}) &=
    \sum_{c,m} w^{(l)}_{c,m}\, \mathcal{N}(\mathbf{x}^{(l)} \mid \boldsymbol{\mu}^{(l)}_{c,m},
    (\sigma^{(l)}_{c,m})^2\mathbf{I}) \\
  p(\mathbf{x}^{(l-1)} \mid \mathbf{x}^{(l)}) &=
    \mathcal{N}\!\left(\mathbf{x}^{(l-1)} \;\middle|\; \Pi_l^\dagger \mathbf{x}^{(l)},\;
    \tfrac{1}{\alpha_{\mathrm{bu}}+\alpha_{\mathrm{td}}}\mathbf{I}\right)
\end{align}

This Bayesian grounding suggests a natural VAE extension, replacing the point-mean-field
approximation with an ELBO objective (Section~\ref{sec:future}).

%% ============================================================
\section{Empirical Validation}
%% ============================================================

We conduct six campaigns across synthetic and real datasets. We report results with
their statistical status (single-seed vs.\ multi-seed, with means $\pm$ standard
deviations where available) and resist drawing conclusions that are not supported by the
sample sizes.

\subsection{Missing Baselines: A Limitation Acknowledged Upfront}
\label{sec:emp:baselines}

We acknowledge that the empirical study does \emph{not} include strong neural baselines
(MLP, CNN) trained under identical protocols. On the datasets we use, established
baselines achieve approximately:
\begin{itemize}
  \item \textbf{MNIST (full)}: logistic regression $\sim 92\%$, MLP $\sim 98\%$,
        CNN $>99\%$.
  \item \textbf{MNIST subset} (1500 train / 500 test, 5 epochs): a 2-layer MLP typically
        reaches $90\%$+ in seconds.
  \item \textbf{Fashion-MNIST (5000 train)}: logistic regression $\sim 84\%$,
        2-layer MLP $\sim 88\%$, small CNN $>90\%$.
\end{itemize}
Our $75.8\%$ on MNIST-subset and $80.1\%$ on Fashion-MNIST are therefore \emph{below
standard baselines}. We do not present HAML as a more accurate classifier on these
datasets. The empirical claims of this preprint concern (i)~the framework trains stably
in high dimension, (ii)~the coupling provides a measurable benefit on structurally
non-convex problems (rings, intermediate-noise moons), and (iii)~the attractor geometry
is interpretable and inspectable. Reaching baseline-competitive accuracy is left as
explicit future work.

\subsection{Campaign A: MNIST Subset --- Gradient Training Gains and Coupling Neutrality}

\textbf{Protocol}: 1500 train / 500 test from MNIST, levels $784 \to 392 \to 196$,
2 attractors per class, 5 epochs, 3-phase training, seed~42.

\begin{table}[h]
\centering
\small
\begin{tabular}{lrr}
\toprule
\textbf{Condition} & \textbf{Test Accuracy} & \textbf{Phase Progress (train)} \\
\midrule
$k$-means init.\ only (no gradient training)       & 45.6\% & --- \\
End of Phase 1 (independent levels)                & ---    & 64.2\% \\
End of Phase 2 (progressive coupling)              & ---    & 68.7\% \\
End of Phase 3 (full coupling)                     & ---    & 71.8\%--77.8\% \\
Full training, independent variant ($\alpha=0$)    & 75.8\% & --- \\
Full training, coupled ($\alpha_{\mathrm{bu}}=\alpha_{\mathrm{td}}=1$) & 75.8\% & --- \\
\bottomrule
\end{tabular}
\caption{Campaign A. Gradient training raises accuracy from $45.6\%$ (k-means
initialisation alone) to $75.8\%$ ($+30.2$~pts). On this short, near-convex protocol,
\textbf{bidirectional coupling provides no measurable benefit} over the independent
variant. Single seed; the $+30.2$~pts gain compares to an extremely weak baseline
($k$-means class assignment with no gradient learning) and should not be interpreted as
evidence of competitiveness with standard MNIST classifiers (see
Section~\ref{sec:emp:baselines}).}
\label{tab:campaign_a}
\end{table}

\textbf{Interpretation.} This is the most important negative result of the study: on a
nearly-convex 10-class problem, coupling is neutral. This is consistent with the theory
(Proposition~\ref{prop:expressivity} requires high topological complexity to predict an
advantage). We treat this result as informative for scoping the regime in which coupling
helps, not as evidence against the architecture.

\subsection{Campaign B: Concentric Rings --- Coupling on Non-Convex Geometry}
\label{sec:emp:rings}

\textbf{Protocol}: alternating concentric rings (topological complexity $\kappa \ge 2$),
$n_{\mathrm{runs}}=5$, seeds~42--46.

\begin{table}[h]
\centering
\small
\begin{tabular}{llr}
\toprule
\textbf{Model} & \textbf{Accuracy (mean $\pm$ std)} & \textbf{Qualitative boundary} \\
\midrule
Independent levels ($\alpha=0$) & $67.30\% \pm 4.37$ & Convex (fails annular topology) \\
Coupled (tuned)                 & $71.20\% \pm 6.43$ & Non-convex (aligns rings) \\
\midrule
$\Delta$ (mean)                 & $+3.90$~pts        & --- \\
\bottomrule
\end{tabular}
\caption{Campaign B. The mean gain ($+3.90$~pts) is smaller than the pooled standard
deviation ($\sigma_{\mathrm{coupled}}=6.43$), so a Welch's two-sample $t$-test on $n=5$
seeds per group does not reject the null hypothesis of equal means at standard
significance levels. However, the qualitative shift in the learned decision boundary
(visible in the per-seed plots in the repository) is consistent across all five coupled
runs.}
\label{tab:campaign_b}
\end{table}

\textbf{Honest reading.} The quantitative mean difference is not statistically significant
at $n=5$ seeds. The qualitative geometric shift (convex decision boundary in the
independent variant; ring-aligned non-convex boundary in the coupled variant) is robust
across all five seeds and is the actual evidence in favour of the coupling on this
benchmark. The high variance in the coupled condition motivates Campaign~E and the
level-aware recovery mechanism.

\subsection{Campaign C: Noisy Moons --- Noise Robustness Profile (Single Seed)}

\textbf{Protocol}: \texttt{make\_moons} with noise $\sigma \in \{0.10, 0.20, 0.30, 0.40\}$,
$n=3000$, \textbf{single seed (42)}.

\begin{table}[h]
\centering
\begin{tabular}{lrrrr}
\toprule
\textbf{Noise} & \textbf{HAML-Indep.} & \textbf{HAML-Coupled} & \textbf{KNN} & \textbf{SVM RBF} \\
\midrule
0.10 & 99.60\% & 99.60\%          & 99.87\% & 99.87\% \\
0.20 & 89.33\% & \textbf{95.60\%} & 96.53\% & 97.20\% \\
0.30 & 88.13\% & 90.53\%          & 90.80\% & 91.33\% \\
0.40 & 84.67\% & 85.07\%          & 83.33\% & 87.33\% \\
\bottomrule
\end{tabular}
\caption{Campaign C. \textbf{Single seed.} The coupling delta
(HAML-Coupled $-$ HAML-Indep.) is $0.00, 6.27, 2.40, 0.40$~pts at noise
$0.10, 0.20, 0.30, 0.40$. The maximum at intermediate noise is consistent with the
prediction that top-down context helps most when local signal is ambiguous but
exploitable. Multi-seed replication is required before treating the bell-shaped profile
as established.}
\label{tab:campaign_c}
\end{table}

\textbf{Falsifiable prediction.} If the bell-shaped profile is real and not a single-seed
artefact, multi-seed runs (10+ seeds) should reproduce: (a)~peak coupling gain in the
intermediate-noise regime, (b)~no significant gain at very low or very high noise. We
state this as a falsifiable prediction the reader can test against the released code.

\subsection{Campaign D: HAML vs.\ a Controlled SA-nODE Reimplementation}

\textbf{Important caveat.} The SA-nODE comparison uses a \emph{controlled internal
reimplementation} of the architecture described in~\cite{marino2024}, not the authors'
official code, which we did not run for this study. The reimplementation reproduces:
single ambient space, binary attractors $\pm a$, analytic double-well potential, linear
coupling layer trained jointly with the dynamics, grid search over $(a, \Delta t, \gamma)$.
A reviewer is justified in regarding this as a weak baseline: a paper comparison should
ultimately use the official code. We report the result as preliminary.

\begin{table}[h]
\centering
\begin{tabular}{lrrr}
\toprule
\textbf{Noise} & \textbf{HAML Coupled} & \textbf{SA-nODE (reimpl.)} & \textbf{Delta} \\
\midrule
0.10 & \textbf{99.60\%} & 86.67\% & $+12.93$~pts \\
0.20 & \textbf{95.60\%} & 86.40\% & $+9.20$~pts \\
0.30 & \textbf{90.53\%} & 85.87\% & $+4.67$~pts \\
0.40 & \textbf{85.07\%} & 83.33\% & $+1.73$~pts \\
\bottomrule
\end{tabular}
\caption{Campaign D. \textbf{Internal SA-nODE reimplementation, single seed.} The
hierarchical architecture dominates our reimplementation of the flat single-space model
across noise levels. SA-nODE degrades more slowly in absolute terms ($-3.33$~pts vs.\
$-14.53$~pts for HAML across the noise range), which we do not currently explain. We
flag this comparison as preliminary pending evaluation against the official SA-nODE code.}
\label{tab:campaign_d}
\end{table}

\subsection{Campaign E: Stabilisation and Inter-Level Collapse}

We iteratively develop the level-aware collapse guard (Section~\ref{sec:p6}) over
subtests E1--E10. Full per-subtest results are in Appendix~A.

\begin{table}[h]
\centering
\small
\begin{tabular}{llrr}
\toprule
\textbf{Subtest} & \textbf{Mechanism} & \textbf{Mean Acc.} & \textbf{Min (worst seed)} \\
\midrule
E8 baseline (coupled)        & None                       & $68.60\% \pm 7.54$ & 56.67\% \\
E8 with level-aware guard    & Sec.~\ref{sec:p6}          & $70.80\% \pm 5.54$ & 60.67\% \\
\midrule
E10 ($n=3200$, seeds 42--46) & Cause-oriented patch       & $78.75\% \pm 7.63$ & 68.25\% \\
\bottomrule
\end{tabular}
\caption{Campaign E summary. The level-aware guard reduces variance and improves the
worst-seed performance on the concentric-rings protocol. Scaling sample size from $1200$
to $3200$ (E10) lifts the mean accuracy by $\approx 8$~pts.}
\label{tab:campaign_e}
\end{table}

\textbf{Mechanistic finding}: collapsed levels show \texttt{separation\_ratio} $\approx 0.15$
(low) despite moderate compactness --- suggesting a learning blockade rather than total
geometric collapse. Re-anchoring of the affected level's attractors resolves the blockade
within 1--2 epochs.

\textbf{Cross-dataset validation} (mutex baseline, seeds 42--46, $n=3200$):
\begin{itemize}
  \item \texttt{noisy\_moons} (noise 0.30): $89.13\% \pm 1.73$, min $86.25\%$.
  \item \texttt{make\_classification} (4 classes): $74.08\% \pm 4.31$, min $67.12\%$.
\end{itemize}

\subsection{Campaign F: Fashion-MNIST --- Stable Training in 784D (Single Seed)}

\textbf{Protocol}: Fashion-MNIST (784D, 10 classes), $n_{\mathrm{train}}=5000$,
$n_{\mathrm{test}}=1000$, levels $784 \to 392 \to 196$, 10 epochs, phases $(2,3,5)$,
\textbf{single seed (42)}.

\begin{table}[h]
\centering
\begin{tabular}{lr}
\toprule
\textbf{Metric} & \textbf{Value} \\
\midrule
Test accuracy                & \textbf{80.1\%} \\
Train accuracy (final)       & 82.3\% \\
End of Phase 1               & 75.36\% \\
End of Phase 2               & 78.86\% \\
End of Phase 3               & 82.26\% \\
Level-aware guard events     & 0 (naturally stable run) \\
Compute (Colab GPU)          & 6\,979~s ($\approx$2~hrs) \\
Reproducibility (Kaggle GPU) & 80.1\% / 6\,222~s \\
\bottomrule
\end{tabular}
\caption{Campaign F. \textbf{Single seed.} The headline result is that the
three-phase training progresses smoothly through all phases in 784D without collapse, not
that $80.1\%$ is a strong accuracy on Fashion-MNIST (it is not; see
Section~\ref{sec:emp:baselines}). Identical results across two GPU environments confirm
deterministic behaviour.}
\label{tab:campaign_f}
\end{table}

\textbf{Compute caveat.} At $\approx$2~hrs per 10-epoch run, a multi-seed Fashion-MNIST
campaign was beyond the compute budget available for this preprint. Multi-seed
validation is the most important empirical follow-up.

%% ============================================================
\section{Mechanistic Diagnostics}
%% ============================================================

We implement diagnostic tools to isolate failure modes:

\subsection{Level-Wise Accuracy Tracking}

At each epoch, per-level accuracy $\mathrm{acc}_l(t)$ is logged. Divergence
$\Delta_l = \max_l \mathrm{acc}_l - \min_l \mathrm{acc}_l$ serves as a collapse indicator.
Typical seed-42 signature on the concentric-rings task:
\begin{itemize}
  \item Phase 1 (independent): all levels converge to $\sim 60\%$.
  \item Phase 2 (progressive): divergence emerges ($\Delta_l \approx 0.15$).
  \item Phase 3 (full coupling): one level plateaus at $\sim 0.5$; guard triggers
        re-anchoring.
\end{itemize}

\subsection{Attractor Geometry Metrics}

Per level:
\begin{itemize}
  \item \texttt{intra\_spread\_mean}: within-class attractor spread (compactness).
  \item \texttt{inter\_centroid\_dist\_min}: minimum between-class centroid distance.
  \item \texttt{separation\_ratio}: geometric margin proxy.
\end{itemize}
Collapsed levels show \texttt{separation\_ratio} $\approx 0.15$ (low) despite non-zero
compactness, indicating a learning blockade rather than total geometric collapse.

\subsection{Force Distribution Analysis}

Histograms of $\|\mathbf{F}^{(l)}\|$ per level reveal:
\begin{itemize}
  \item \emph{Saturated kernels} (pre-fix~I1 in high dimensions): $\|\mathbf{F}\| \approx 0$.
  \item \emph{Balanced learning} (post-fix~I1): approximately Gaussian distribution.
  \item \emph{Gradient death} in late phases: indicator of convergence or instability.
\end{itemize}

%% ============================================================
\section{Discussion}
%% ============================================================

\subsection{What the Evidence Supports, and What It Does Not}

\textbf{Evidence supports:}
\begin{enumerate}
  \item The framework trains stably in dimensions up to 784, given condition~I1.
  \item On structurally non-convex problems (concentric rings), coupled HAML
        \emph{qualitatively} recovers ring-aligned decision boundaries that the
        independent variant fails to produce.
  \item Bidirectional coupling provides a single-seed accuracy gain at intermediate
        noise on noisy moons (consistent with theory; awaiting multi-seed confirmation).
  \item Exact MAP correspondence with hierarchical Predictive Coding under explicit
        conditions M1--M3 (as a specialisation of the general PC--inference
        correspondence; see Remark~\ref{rem:map_priorwork}).
\end{enumerate}

\textbf{Evidence does not support:}
\begin{enumerate}
  \item HAML being competitive with standard neural baselines on MNIST or Fashion-MNIST
        on accuracy alone.
  \item A statistically significant mean accuracy gain from coupling at $n=5$ seeds on
        rings.
  \item Superiority over SA-nODE based on the authors' official implementation.
  \item Scalability beyond 784D.
  \item Novelty of any individual building block taken in isolation.
\end{enumerate}

\subsection{Genuine Limitations}
\label{sec:discussion:limitations}

\begin{enumerate}
  \item \textbf{Combinatorial rather than foundational novelty}: each architectural
        primitive of HAML --- continuous attractors, hierarchical PC, energy
        minimisation, fixed-point computation, dimension-scaled kernels --- has direct
        antecedents in the literature reviewed in Section~\ref{sec:related}. The
        contribution is the specific assembly of these primitives for supervised
        classification by basin membership, not the invention of any of them. Readers
        should evaluate HAML as a synthesis with motivated specialisation, not as a
        new family of models.
  \item \textbf{Missing strong baselines}: no MLP, CNN, or boosted-tree baseline trained
        under identical protocols. This is the most important methodological gap.
  \item \textbf{Single-seed runs on real data}: MNIST-subset and Fashion-MNIST are
        single-seed for compute reasons.
  \item \textbf{Statistical underpowering on rings}: $n=5$ seeds is insufficient to claim
        a mean accuracy advantage from coupling; we rely on qualitative geometric
        evidence.
  \item \textbf{Internal SA-nODE reimplementation}: not a substitute for evaluation
        against the authors' official code.
  \item \textbf{Computational cost}: $\approx$2~hrs per 10-epoch run on Fashion-MNIST
        (5000 samples). Adjoint methods and gradient checkpointing should reduce this
        substantially.
  \item \textbf{Inter-level instability}: variance remains high between seeds; the
        level-aware guard helps but does not fully resolve it.
  \item \textbf{Hyperparameter sensitivity}: $(\alpha_{\mathrm{bu}}, \alpha_{\mathrm{td}}, \mu_1, \mu_2)$
        require per-dataset tuning. Automated search is an open engineering challenge.
  \item \textbf{Unexplained asymmetry}: SA-nODE-reimpl.\ degrades more slowly than HAML
        as noise increases (Campaign~D); we do not currently have a mechanistic
        explanation.
  \item \textbf{Conjecture status}: Proposition~\ref{prop:bidi_necessary} (bidirectional
        necessity) is stated as a conjecture without a proof.
  \item \textbf{Theoretical specialisation, not foundational novelty}: the theoretical
        results in Section~\ref{sec:theory} are specialisations of existing frameworks
        (LaSalle stability for gradient flows; Rademacher complexity for neural ODEs
        \cite{hanson2024rademacher,marion2023generalization}; PC--inference
        correspondence \cite{whittington2017approximation,song2020can,salvatori2021exactbp})
        to HAML's specific architecture, not foundational results in their own right.
\end{enumerate}

\subsection{Future Work}
\label{sec:future}

\begin{enumerate}
  \item \textbf{Strong baselines}: MLP / CNN / boosted trees under identical protocols on
        all reported datasets.
  \item \textbf{Multi-seed Fashion-MNIST}: at least 5 seeds with the adjoint method to
        amortise compute.
  \item \textbf{Official SA-nODE comparison}: replication against Marino et al.'s released
        code.
  \item \textbf{Adjoint integration}: full neural-ODE adjoint sensitivity to reduce
        memory and training time.
  \item \textbf{Generative extension (HAML-VAE)}: replace the point-mean-field
        approximation with an ELBO objective.
  \item \textbf{Proof of Conjecture~\ref{prop:bidi_necessary}}: a constructive separation
        between bidirectional and independent-levels representations.
  \item \textbf{Tighter Rademacher constants} and sample complexity bounds.
  \item \textbf{Scalability beyond 784D}: tested on CIFAR-10 (3072D) and larger.
\end{enumerate}

%% ============================================================
\section{Reproducibility and Implementation}
%% ============================================================

All code is available at:
\url{https://github.com/ikobootloader/HIERARCHICAL-ATTRACTOR-MACHINE-LEARNING}

\subsection{Key Software}

\begin{itemize}
  \item \texttt{haml/}: main package (dynamics, model, training, projection, analysis,
        visualisation).
  \item \texttt{experiments/}: 6 reproducible campaign scripts (A--F).
  \item \texttt{diagnose\_forces.py}: high-dimension saturation detector.
  \item Dependencies: \texttt{numpy, scikit-learn, torch, torchdiffeq, scipy}.
\end{itemize}

\subsection{Reproducible Commands}

\textbf{Campaign~A} (MNIST ablation):
\begin{verbatim}
python experiments/mnist_coupling_ablation.py
\end{verbatim}

\textbf{Campaign~B} (concentric rings, multi-seed):
\begin{verbatim}
python experiments/concentric_coupling_ablation.py --n-runs 5 --save-figure
\end{verbatim}

\textbf{Campaign~C} (noisy benchmark):
\begin{verbatim}
python experiments/noisy_benchmark.py
\end{verbatim}

\textbf{Campaign~F} (Fashion-MNIST, GPU):
\begin{verbatim}
PYTHONIOENCODING=utf-8 python experiments/fashion_mnist_short_probe.py \
  --device cuda --seed 42 --n-train 5000 --n-test 1000 --n-epochs 10 \
  --phase1-epochs 2 --phase2-epochs 3
\end{verbatim}

%% ============================================================
\section{Conclusion}
%% ============================================================

We introduced HAML, a hierarchical bidirectional attractor framework for supervised
classification, combining PCA-decomposed state spaces, learned continuous attractors, and
Predictive-Coding-style coupling. The core architectural contribution is the integration
of these three elements within a single dissipative dynamical system, together with a
practical training recipe that empirically stabilises it.

Theoretical contributions comprise: dissipative convergence (Proposition~\ref{prop:convergence}),
dimension-independent local rate at fixed points (Proposition~\ref{prop:dimindep}),
a hierarchical expressivity argument (Propositions~\ref{prop:expressivity}--\ref{prop:rademacher}),
and an exact correspondence with hierarchical Predictive Coding MAP inference under
explicit conditions (Theorem~\ref{thm:map}).

Empirically, the evidence is preliminary and mixed. The framework trains stably in 784D
($80.1\%$ Fashion-MNIST, single seed, below standard baselines). Coupling provides a
qualitative geometric advantage on structurally non-convex problems (rings) and a
single-seed accuracy gain at intermediate noise (moons), but no measurable gain on
near-convex problems (MNIST subset). Multi-seed validation and comparison against strong
neural baselines remain the most important next steps.

We release the framework and all experimental code in the hope that it will be useful as
a starting point for further investigation of attractor-based classification, hierarchical
Predictive Coding implementations, and energy-based hybrids.

\section*{Acknowledgments}

The author thanks the arXiv and open-source scientific communities, and the contributors
to PyTorch, scikit-learn, SciPy, and \texttt{torchdiffeq}. This project originated after
reading J.~Gribbin, \textit{Deep Simplicity} (2004).

%% ============================================================
\begin{thebibliography}{99}

\bibitem{amit2001attractor}
Amit, Y., \& Mascaro, M. (2001).
Attractor networks for shape recognition.
\textit{Neural Computation}, \textbf{13}(6), 1415--1442.
DOI:10.1162/08997660152002906.

\bibitem{bai2019deep}
Bai, S., Kolter, J. Z., \& Koltun, V. (2019).
Deep equilibrium models.
\textit{Advances in Neural Information Processing Systems (NeurIPS)}, 32.
arXiv:1909.01377.

\bibitem{bai2020multiscale}
Bai, S., Koltun, V., \& Kolter, J. Z. (2020).
Multiscale deep equilibrium models.
\textit{Advances in Neural Information Processing Systems (NeurIPS)}, 33, 5238--5250.
arXiv:2006.08656.

\bibitem{chen2018neural}
Chen, T. Q., Rubanova, Y., Bettencourt, J., \& Duvenaud, D. K. (2018).
Neural ordinary differential equations.
\textit{Advances in Neural Information Processing Systems}, 31, 6571--6583.

\bibitem{friston2005theory}
Friston, K. (2005).
A theory of cortical responses.
\textit{Philosophical Transactions of the Royal Society B: Biological Sciences},
\textbf{360}(1456), 815--836.

\bibitem{friston2017active}
Friston, K. (2017).
Active inference: A process theory.
\textit{Neural Computation}, \textbf{29}(1), 1--49.

\bibitem{han2018deep}
Han, K., Wen, H., Zhang, Y., Fu, D., Culurciello, E., \& Liu, Z. (2018).
Deep predictive coding network with local recurrent processing for object recognition.
\textit{Advances in Neural Information Processing Systems (NeurIPS)}, 31, 9221--9233.
arXiv:1805.07526.

\bibitem{hanson2024rademacher}
Hanson, J., \& Raginsky, M. (2024).
Rademacher complexity of neural ODEs via Chen--Fliess series.
\textit{Proceedings of the 6th Annual Learning for Dynamics \& Control Conference (L4DC)},
PMLR 242:758--769.
arXiv:2401.16655.

\bibitem{krotov2016dense}
Krotov, D., \& Hopfield, J. J. (2016).
Dense associative memory for pattern recognition.
\textit{Advances in Neural Information Processing Systems}, 29.
arXiv:1606.01164.

\bibitem{krotov2021hierarchical}
Krotov, D. (2021).
Hierarchical associative memory.
\textit{arXiv preprint} arXiv:2107.06446.

\bibitem{larochelle2008}
Larochelle, H., \& Bengio, Y. (2008).
Classification using discriminative restricted Boltzmann machines.
\textit{Proceedings of the 25th International Conference on Machine Learning (ICML)},
536--543.

\bibitem{marino2024}
Marino, R., Giambagli, L., Chicchi, L., Buffoni, L., \& Fanelli, D. (2024).
Stable attractors for neural networks classification via ordinary differential equations
(SA-nODE).
\textit{arXiv preprint} arXiv:2311.10387v2.

\bibitem{marion2023generalization}
Marion, P. (2023).
Generalization bounds for neural ordinary differential equations and deep residual
networks.
\textit{Advances in Neural Information Processing Systems (NeurIPS)}, 36.

\bibitem{rao1999predictive}
Rao, R. P. N., \& Ballard, D. H. (1999).
Predictive coding in the visual cortex: A functional interpretation of some
extra-classical receptive-field effects.
\textit{Nature Neuroscience}, \textbf{2}(1), 79--87.

\bibitem{ramsauer2021hopfield}
Ramsauer, H., Sch\"{a}fl, B., Lehner, J., Seidl, P., Widrich, M., Gruber, L.,
\ldots\ \& Kreil, D. P. (2021).
Hopfield networks is all you need.
\textit{International Conference on Learning Representations (ICLR)}.

\bibitem{salvatori2021associative}
Salvatori, T., Song, Y., Hong, Y., Frieder, S., Sha, L., Xu, Z., Bogacz, R.,
\& Lukasiewicz, T. (2021).
Associative memories via predictive coding.
\textit{Advances in Neural Information Processing Systems (NeurIPS)}, 34.
arXiv:2109.08063.

\bibitem{salvatori2021exactbp}
Salvatori, T., Song, Y., Lukasiewicz, T., Bogacz, R., \& Xu, Z. (2021).
Predictive coding can do exact backpropagation on any neural network.
\textit{arXiv preprint} arXiv:2103.04689.

\bibitem{sato1996glvq}
Sato, A., \& Yamada, K. (1996).
Generalized learning vector quantization.
\textit{Advances in Neural Information Processing Systems}, 8.

\bibitem{song2020can}
Song, Y., Lukasiewicz, T., Xu, Z., \& Bogacz, R. (2020).
Can the brain do backpropagation? --- Exact implementation of backpropagation in
predictive coding networks.
\textit{Advances in Neural Information Processing Systems (NeurIPS)}, 33, 22566--22579.

\bibitem{spratling2017hierarchical}
Spratling, M. W. (2017).
A hierarchical predictive coding model of object recognition in natural images.
\textit{Cognitive Computation}, \textbf{9}(2), 151--167.
DOI:10.1007/s12559-016-9445-1.

\bibitem{weyl1912}
Weyl, H. (1912).
Das asymptotische Verteilungsgesetz der Eigenwerte linearer partieller
Differentialgleichungen.
\textit{Mathematische Annalen}, \textbf{71}, 441--479.

\bibitem{whittington2017approximation}
Whittington, J. C. R., \& Bogacz, R. (2017).
An approximation of the error backpropagation algorithm in a predictive coding network
with local Hebbian synaptic plasticity.
\textit{Neural Computation}, \textbf{29}(5), 1229--1262.
DOI:10.1162/NECO\_a\_00949.

\bibitem{winston2020monotone}
Winston, E., \& Kolter, J. Z. (2020).
Monotone operator equilibrium networks.
\textit{Advances in Neural Information Processing Systems (NeurIPS)}, 33.
arXiv:2006.08591.

\end{thebibliography}

%% ============================================================
\appendix
%% ============================================================

\section{Detailed Campaign Results}

\subsection{Campaign~E --- Subtest Summary}

\begin{table}[h]
\centering
\footnotesize
\setlength{\tabcolsep}{4pt}
\begin{tabular}{llrr}
\toprule
\textbf{Subtest} & \textbf{Key mechanism} & \textbf{Accuracy (test)} & \textbf{$\Delta$ vs.\ indep.} \\
\midrule
E1  & Reinforced coupling ($\alpha_{\mathrm{td}}=1.5$, 25 ep.) & 86.38\%            & $+15.63$ pts \\
E2  & Early-stop + LR decay                                    & 82.88\%            & $+12.00$ pts \\
E3  & Adaptive $\mu_{\mathrm{sep}}$ + soft landing ep.\ 20     & 69.88\%            & $-1.25$ pts  \\
E4  & Soft landing on divergence trigger                       & 76.25\%            & $+6.37$ pts  \\
E5  & Raised trigger + patience 2                              & 76.13\%            & $+5.00$ pts  \\
E6  & Collapse guard on divergence jump                        & 86.38\%            & $+16.00$ pts \\
E7  & Seed-wise diagnostics + LR cascade patch                 & $69.33\% \pm 7.10$ & $+9.60$ pts  \\
E8  & Level-aware guard (phase 3 only)                         & $70.80\% \pm 5.54$ & $+11.07$ pts \\
E10 & Scale $n=3200$, cause-oriented patch                     & $78.75\% \pm 7.63$ & ---          \\
\bottomrule
\end{tabular}
\caption{Stabilisation iteration summary (Campaign~E). Results on the concentric-rings
protocol, except E10 where $n=3200$, seeds~42--46. The high single-seed peaks (E1, E6)
were not retained because their stability across seeds is not established. E8 was
selected as the operational mechanism on the multi-seed mean--variance trade-off.}
\label{tab:campaign_e_detail}
\end{table}

\subsection{Initialisation Condition I1 --- Justification}

Distance concentration in $\mathbb{R}^d$:
\begin{equation}
  \mathbb{E}[\|\mathbf{x} - \mathbf{y}\|] \approx \sqrt{d}\cdot\sigma_{\mathrm{data}}
\end{equation}
For a Gaussian kernel $\exp(-r^2/2\sigma^2)$ to maintain activations in the range
$[0.1, 0.9]$ across dimension $d$, $\sigma$ must scale as $\sqrt{d}$. Without~I1, the
following table shows the resulting kernel saturation (single typical run).

\begin{table}[h]
\centering
\begin{tabular}{rrrr}
\toprule
\textbf{Dimension} & \textbf{Pre-fix kernel} & \textbf{Post-fix kernel} & \textbf{Gradient} \\
\midrule
10  & 0.95    & 0.88 & 0.24 \\
100 & 0.03    & 0.67 & 0.18 \\
784 & $<0.01$ & 0.52 & 0.15 \\
\bottomrule
\end{tabular}
\caption{Kernel activations and gradient magnitudes with and without condition~I1.}
\end{table}

\section{Hyperparameter Sensitivity}

Ablation on the concentric-rings protocol:

\begin{table}[h]
\centering
\begin{tabular}{lrr}
\toprule
\textbf{Parameter} & \textbf{Range tested} & \textbf{Accuracy range} \\
\midrule
$\alpha_{\mathrm{bu}}$ & $[0.0, 1.0]$   & $[67.3\%, 73.2\%]$ \\
$\alpha_{\mathrm{td}}$ & $[0.0, 1.5]$   & $[67.3\%, 75.1\%]$ \\
$\mu_1$ (separation)   & $[0.1, 1.0]$   & $[70.5\%, 71.9\%]$ \\
$\mu_2$ (dynamics)     & $[0.001, 0.1]$ & $[71.1\%, 71.3\%]$ \\
\bottomrule
\end{tabular}
\caption{Sensitivity on primary hyperparameters. The model is moderately robust;
$\alpha_{\mathrm{bu}}$ and $\alpha_{\mathrm{td}}$ require the most attention per dataset.}
\end{table}

\section{Summary of Theoretical Results}

\begin{table}[h]
\centering
\begin{tabular}{llp{6.5cm}}
\toprule
\textbf{Result} & \textbf{Section} & \textbf{Informal statement} \\
\midrule
Prop.~\ref{prop:convergence} (Convergence)        & \ref{sec:theory:convergence} & Bounded trajectories converge to critical points; no divergence or sustained oscillation. \\
Prop.~\ref{prop:correct_basin} (Correct basin)    & \ref{sec:theory:convergence} & Under C1--C3, gradient flow stays in the correct basin (conditional, not deep). \\
Prop.~\ref{prop:guidance} (Hierarchical guide)    & \ref{sec:theory:convergence} & Eigenvalue analysis shows coupling accelerates local convergence. \\
Prop.~\ref{prop:dimindep} (Dim.\ independence)    & \ref{sec:theory:dimindep}    & Local rate at a fixed point is dimension-independent. \\
Prop.~\ref{prop:expressivity} (Expressivity)      & \ref{sec:theory:expressivity} & Exponential parameter gain on hierarchically factorisable problems. \\
Conj.~\ref{prop:bidi_necessary} (Bidi.\ necessity) & \ref{sec:theory:expressivity} & Bidirectional coupling is conjectured necessary on certain problems (no proof). \\
Prop.~\ref{prop:rademacher} (Rademacher)          & \ref{sec:theory:expressivity} & Sub-linear-in-depth complexity bound. \\
Thm.~\ref{thm:map} (MAP equivalence)              & \ref{sec:theory:bayes}       & Exact MAP in hierarchical Gaussian-mixture model under M1--M3. \\
Prop.~\ref{prop:p6} (Level-aware recovery)         & \ref{sec:p6}                 & Heuristic engineering mechanism reduces inter-seed variance. \\
\bottomrule
\end{tabular}
\caption{Summary of theoretical results, with their epistemic status.}
\end{table}

\end{document}
